\section*{अध्याय \ref{ch:b3:deformation-mechanisms} --- विरूपण-क्रियाविधि और अपरूपण क्षेत्र}

\begin{solution}{exo:b3:deformation-mechanisms:1}
फ़ॉर्स्टेराइट: $1573/2163 = 0.73$; बर्फ़: $263/273 = 0.96$। बर्फ़ अपने गलनांक के
बहुत निकट है, इसलिए अधिक सरलता से विसर्पित होती है: हिमनद पृथ्वी की
सतह पर प्रवाहित होते हैं।
\end{solution}

\begin{solution}{exo:b3:deformation-mechanisms:2}
$\gamma = \tan\ang{70} = 2.75$ और $\tan\ang{80} = 5.67$। \ang{90} के निकट
स्पर्शज्या कोण के साथ बहुत तेज़ी से बदलती है: एक अंश की त्रुटि $\psi =
\ang{84}$ पर $\gamma$ को लगभग एक-पाँचवाँ भाग बदल देती है।
\end{solution}

\begin{solution}{exo:b3:deformation-mechanisms:3}
$\sigma = (10^{3.56}/10)^{1/1.26} = 363.1^{0.794} = \qty{108}{MPa}$।
\end{solution}

\begin{solution}{exo:b3:deformation-mechanisms:4}
$E/RT = 223000/(8.314\times773.15) = 34.69$ है, अतः $\dot\varepsilon =
\num{1.1e-4}\times20^4\times e^{-34.69} = \num{1.1e-4}\times\num{1.6e5}\times
\num{8.6e-16} = \qty{1.5e-14}{s^{-1}}$।
\end{solution}

\begin{solution}{exo:b3:deformation-mechanisms:5}
$\eta = \sigma/(2\dot\varepsilon)$: $\num{65.4e6}/\num{2e-14} =
\qty{3.3e21}{Pa.s}$ पर \qty{400}{\celsius} और $\num{18.0e6}/\num{2e-14} =
\qty{9.0e20}{Pa.s}$ पर \qty{500}{\celsius}: इन तापमानों पर मध्य भूपर्पटी
लगभग उतनी ही श्यान है जितना हिमचादरों के नीचे का प्रावार।
\end{solution}

\begin{solution}{exo:b3:deformation-mechanisms:6}
$n = \log(\num{1.6e-4}/\num{e-5})/\log2 = \log16/\log2 = 4$: \omterm{def:b3:deformation-mechanisms:creep}{अव्यवस्थापन विसर्पण},
जिसका \omterm{def:b3:deformation-mechanisms:flow-law}{प्रतिबल घातांक} 1 से कहीं अधिक है।
\end{solution}

\begin{solution}{exo:b3:deformation-mechanisms:7}
$\ln10 = (E/R)(1/1000 - 1/1100)$, अतः $E = 8.314\times2.3026/\num{9.091e-5} =
\qty{211}{kJ.mol^{-1}}$।
\end{solution}

\begin{solution}{exo:b3:deformation-mechanisms:8}
नैबैरो--हेरिंग ($d^{-2}$): 100 गुना तेज़; कोबल ($d^{-3}$): 1000 गुना।
सीमाओं के साथ विसरित होने वाले परमाणु ऐसा पथ अपनाते हैं जिसका अनुप्रस्थ काट, $\delta
d$, कण के साथ घटता है; इसके अतिरिक्त छोटी दूरी और
स्थानांतरित करने के लिए कम आयतन भी काम आता है।
\end{solution}

\begin{solution}{exo:b3:deformation-mechanisms:9}
$\sigma_b^{2.5} = \num{3.59e-6}\,d^{-3}/\num{6.54e-17}$: $d = 1000$ के लिए,
$\sigma_b = 54.9^{0.4} = \qty{5.0}{MPa}$; $d = 100$ के लिए, $\sigma_b =
\num{54900}^{0.4} = \qty{79}{MPa}$। ढलान $-3/2.5 = -1.2$ है: दस गुना
बारीक कण सीमा-प्रतिबल को $10^{1.2} = 16$ गुना बढ़ा देता है।
\end{solution}

\begin{solution}{exo:b3:deformation-mechanisms:10}
दोनों पथों में सांद्रता का प्रेरक अंतर समान है और
दूरी $d$ है; प्रति इकाई विकृति में स्थानांतरित होने वाला आयतन $d^3$ है।
जालक का अनुप्रस्थ काट $\propto d^2$ है, अतः दर $d^2/(d\,d^3) =
d^{-2}$ है; सीमाएँ $\propto\delta d$ देती हैं, अतः $\delta d/(d\,d^3) = \delta
d^{-3}$। सीमा और जालक की दरों का अनुपात $(B_B\delta D_B)/(B_LD_Ld)$ है,
1 से बड़ा होता है, जब $d < (B_B/B_L)\,\delta D_B/D_L$।
\end{solution}

\begin{solution}{exo:b3:deformation-mechanisms:11}
$\eta \propto \dot\varepsilon^{(1-n)/n} = \dot\varepsilon^{-2.5/3.5}$: दस गुना
तेज़ होने पर श्यानता $10^{0.714} = 5.2$ गुना घट जाती है। जो पट्टी संयोगवश
तेज़ विरूपित होती है, वह दुर्बल हो जाती है, इसलिए विरूपण का अधिक भाग वही लेती है:
विकृति उसी में एकत्र हो जाती है।
\end{solution}

\begin{solution}{exo:b3:deformation-mechanisms:12}
पहले: $\dot\varepsilon_{\mathrm{dis}} = \num{6.54e-17}\times10^{3.5} =
\qty{2.07e-13}{s^{-1}}$ और $\dot\varepsilon_{\mathrm{dif}} = \num{3.59e-6}
\times10/3000^3 = \qty{1.3e-15}{s^{-1}}$: \qty{2.08e-13}{s^{-1}}, \omterm{def:b3:deformation-mechanisms:creep}{अव्यवस्थापन
विसर्पण}। बाद में: $\dot\varepsilon_{\mathrm{dif}} = \num{3.59e-6}\times10/30^3 =
\qty{1.33e-9}{s^{-1}}$, \omterm{def:b3:deformation-mechanisms:creep}{विसरण विसर्पण}, वही प्रतिबल पर \num{6400} गुना तेज़:
पुनःक्रिस्टलन ने क्षेत्र को बहुत अधिक मृदु कर दिया है।
\end{solution}

\begin{solution}{pb:b3:deformation-mechanisms:1}
\textbf{1.} क्षेत्र के सापेक्ष डाइक का कोण \ang{90} से घटकर \ang{6} हो जाता है: $\psi =
\ang{84}$ और $\gamma = \tan\ang{84} = 9.51$।

\textbf{2.} $\gamma w = 9.51\times400 = \qty{3804}{m}$, लगभग \qty{3.8}{km}।

\textbf{3.} $S_1 = \sqrt{1 + 9.51^2/4} + 9.51/2 = 4.859 + 4.755 = 9.614$,
$S_2 = 1/S_1 = 0.104$, और $S_1/S_2 = 92.4$।

\textbf{4.} $\tan2\theta' = 2/9.51 = 0.210$, $\theta' = \ang{5.9}$: पर्णन,
दीर्घ अक्ष के निकट रहकर, क्षेत्र की ओर जैसे-जैसे $\gamma$ बढ़ता है,
और अधिक निकट मुड़ता जाता है।

\textbf{5.} उत्तर की ओर अनुसरण करने पर डाइक क्षेत्र के भीतर पूर्व की ओर ले जाई जाती है और
उससे आगे पूर्व की ओर निकलती है: उत्तरी भित्ति दक्षिणी के सापेक्ष पूर्व की ओर खिसकी।
दक्षिणी भित्ति से देखने पर दूर का पक्ष दाईं ओर खिसका है: क्षेत्र
दाएँ-पार्श्वीय है।

\textbf{6.} खिंचाव-रेखन वाले तल में, जो स्ट्राइक-स्लिप क्षेत्र के लिए
यहाँ क्षैतिज है, और ऊर्ध्वाधर पर्णन के लंबवत: एक क्षैतिज काट।
इसका S पर्णन क्षेत्र से वामावर्त झुकता है,
पूर्व-उत्तर-पूर्व की ओर, विकृति-दीर्घवृत्त के दीर्घ अक्ष की भाँति।

\textbf{7.} $\sigma = (10^{3.56}/20)^{1/1.26} = 181.5^{0.794} = \qty{62}{MPa}$।

\textbf{8.} $242.1^{0.794} = \qty{78}{MPa}$ के लिए \qty{15}{\micro m} और
$145.2^{0.794} = \qty{52}{MPa}$ के लिए \qty{25}{\micro m}।

\textbf{9.} हाँ: प्रयोग 34 से \qty{268}{MPa} तक के परास में फैले थे।

\textbf{10.} यह क्षेत्र \qty{450}{\celsius} पर विरूपित हुआ, जो प्रयोगों के 800 से
\qty{1100}{\celsius} तक के तापमान से बहुत नीचे था; तापमान पर निर्भर न करने वाला
संबंध इस तापमान तक ले जाया जा सकता है, यद्यपि यह अब भी एक
बहिर्वेशन है।

\textbf{11.} $D = 3631/100^{1.26} = 3631/331 = \qty{11}{\micro m}$।

\textbf{12.} $E/RT = 223000/(8.314\times723.15) = 37.09$ और $e^{-37.09} =
\num{7.79e-17}$।

\textbf{13.} $\dot\varepsilon = \num{1.1e-4}\times62^4\times\num{7.79e-17} =
\num{1.1e-4}\times\num{1.478e7}\times\num{7.79e-17} = \qty{1.27e-13}{s^{-1}}$।

\textbf{14.} $\dot\gamma = 2\dot\varepsilon = \qty{2.53e-13}{s^{-1}}$।

\textbf{15.} $\eta = \num{62e6}/(2\times\num{1.27e-13}) = \qty{2.4e20}{Pa.s}$,
\qty{3e20}{Pa.s} के निकट है, जो आर्डो झील के नीचे के ऊपरी प्रावार का है;
यह क्षेत्र प्रावार जितना ही दुर्बल था।

\textbf{16.} $e^{-E/RT}$ \qty{400}{\celsius} पर गुणक 0.064 से
बदलता है (सोलह गुना धीमा), और \qty{500}{\celsius} पर 11.0 गुणक से।

\textbf{17.} $(78/62)^4 = 2.5$ गुना तेज़।

\textbf{18.} $t = 9.51/\num{2.53e-13} = \qty{3.75e13}{s}$।

\textbf{19.} $\num{3.75e13}/\num{3.156e10} = \qty{1190}{kyr}$।

\textbf{20.} $\qty{3804}{m}/\qty{1.19e6}{yr} = \qty{3.2}{mm/yr}$।

\textbf{21.} मुख्य प्लेट सीमाओं की तुलना में दस या अधिक गुना धीमा:
यह क्षेत्र एक मामूली शाखा था, या कई शाखाओं में से एक, जो किसी सीमा की गति साझा करती थीं।

\textbf{22.} $\dot\gamma = 0.02/(\num{3.156e7}\times400) = \qty{1.58e-12}{s^{-1}}$,
6.2 गुना तेज़; प्रतिबल $6.2^{1/4} = 1.58$ गुना बढ़कर \qty{98}{MPa} हो जाता है, और
कणों का आकार $3631/98^{1.26} = \qty{11}{\micro m}$ होगा।

\textbf{23.} समय $\sigma^{-4}$ के अनुपात में बदलता है: केवल कण-आकार का परास
$(62/78)^4 = 0.40$ से $(62/52)^4 = 2.0$ गुना तक अनुमान देता है, और
घातांक $1.26 \pm 0.13$ और \omterm{def:b3:deformation-mechanisms:flow-law}{प्रवाह नियम} के स्थिरांकों की अनिश्चितता
इसे और चौड़ा कर देती है: परिणाम एक कोटि का अनुमान है, कोई तिथि नहीं।

\textbf{24.} वार्डन \omterm{def:b3:deformation-mechanisms:shear-zone}{अपरूपण क्षेत्र} को अपनी विकृति संचित करने में लगभग
\textbf{\qty[text-series-to-math=true]{1190}{kyr}} लगा।
\end{solution}
